% ==========================================================
% titles/title-10.tex -- Title X: Reparative Justice and
% Indigenous Sovereignty
% ==========================================================

\ConstTitle{Reparative Justice and Indigenous Sovereignty}{title:reparative}

\begin{quote}
    \itshape
    This Commonwealth acknowledges that any territory in which it is established will carry a history of dispossession of indigenous peoples, and of accumulated wealth and social structures built through the enslaved and coerced labor of peoples subjected to racial domination. These historical harms produce ongoing inequalities that universal systems alone do not correct.
\end{quote}

\Article{Indigenous Sovereignty}{art:indigenous-sov}

\ArtSection{The Land Back Negotiation Office}{art:indigenous-sov:s-land-back-office}
A \gls{gls:land-back-negotiation-office}, jointly structured by the Commonwealth and indigenous nations with territorial claims on Commonwealth land, is established within one year of ratification. Its mandate, staffing, and budget are co-determined in a preliminary convening achieving representation from all indigenous nations with territorial claims. Staffing term structures are determined through the joint convening process. If the convening has not reached agreement on staffing, budget, or term structures by the one-year establishment deadline, the \gls{gls:electoral-commons} convenes an interim Office at a scale sufficient to carry on negotiations, exercising no discretion over the outcome and making no substantive negotiating decisions; this interim arrangement is without prejudice to the permanent structure the convening ultimately adopts.

\ArtSection{Negotiation Timeline and Review}{art:indigenous-sov:s-negotiation-timeline}
A jointly agreed negotiation framework is published within two years. Annual good faith progress reports are published jointly beginning in year three. A mandatory \gls{gls:constitutional-assembly} review occurs in year seven, with remedial authority for the \gls{gls:constitutional-court} if the Commonwealth has failed to negotiate in good faith.

\ArtSection{Sovereignty Exit}{art:indigenous-sov:s-sovereignty-exit}
Where an indigenous nation chooses to exercise full sovereignty over returned land, that land exits Commonwealth territory entirely. The relationship between the Commonwealth and sovereign indigenous nations is thereafter governed by Article \ref{art:foreign-relations} and Article \ref{art:anti-imperialism}; the anti-imperialism principle applies in full and the Commonwealth may not use economic leverage, infrastructure access, or resource dependency as instruments of coercion toward sovereign indigenous nations.

\ArtSection{Transition Protections}{art:indigenous-sov:s-transition-protections}
The terms of infrastructure transition on returned land are determined through negotiation; the Commonwealth may not use infrastructure dependency as leverage in those negotiations. The rights of Commonwealth residents present on returned land during any sovereignty transition are addressed within the negotiated agreement; the Commonwealth commits to protecting those residents' Material Sufficiency Floor rights (Article \ref{art:matsuff}) during any transition period, without using their presence as grounds to delay or deny sovereignty recognition.

\ArtSection{Entrenchment}{art:indigenous-sov:s-entrenchment}
Territorial agreements reached through this process are constitutionally entrenched upon ratification by both parties and may not be reversed by subsequent legislative majorities on grounds of territorial integrity. The Commonwealth's territorial boundary is subject to justice-based revision through this process; this commitment is permanent and may not be abolished by any \gls{gls:constitutional-assembly}.

\Article{Reparations}{art:reparations}

\ArtSection{The Reparations Design Council}{art:reparations:s-council}
A \gls{gls:reparations-design-council} with affected community majority representation and genuine veto authority over its mandate is constituted within one year of ratification. Initial members are selected through nomination processes organized by affected communities and coordinated administratively by the \gls{gls:electoral-commons}, which verifies eligibility but exercises no discretion over who is nominated. Once constituted, the Council determines its own ongoing term structure and future selection process, consistent with its mandate and affected community majority governance. Land-based remedies are within the Council's authority to recommend.

\ArtSection{Timeline and Funding Protection}{art:reparations:s-timeline-funding}
A documented Harm Assessment is published within two years. A Form and Scale Recommendation is published within three years; the \gls{gls:sortition-legislature} may not reduce the recommended scale by more than twenty percent without supermajority vote and \gls{gls:constitutional-court} review for consistency with the anti-domination principle. Annual reparative allocations begin in year four, protected from reduction below the recommended scale by the same supermajority threshold.

\ArtSection{Review and Completion Standard}{art:reparations:s-completion-standard}
A mandatory ten-year review assesses gap closure, with automatic escalation if gaps are not closing. The reparative obligation continues until the \gls{gls:reparations-design-council}, with an affected-community majority, formally determines it has been sufficiently addressed. The Commonwealth does not unilaterally declare this obligation complete.

\ArtSection{Calibration to Territorial History}{art:reparations:s-calibration}
The specific forms of harm assessment and remedy are calibrated to the territory's actual history; the \gls{gls:reparations-design-council}'s first task is to determine which historical harms apply in the specific territorial context of the Commonwealth's establishment.
